\section{Lean formalization and ProveIt integration}\label{sec:lean}

\subsection{Verification status}

The archive deliberately separates three kinds of evidence.
The manuscript contains the complete human-readable proof
and its iPRA induction audit. The executable code checks
the concrete finite algorithm on specified instances.
The file \code{lean/FiniteCore.lean} contains compiled
host-Lean proofs of finite arithmetic and algebra.
The complete implicit function theorem has not been
formalized in that file, and the archive does not contain
a completed \(\Prf\,\PRAAx\) derivation of the full theorem.

The Lean target imports only the bundled \code{Std}
library and was checked with Lean 4.32.0, compiler commit
\code{8c9756b28d64dab099da31a4c09229a9e6a2ef35}.
Both direct compilation and \code{lake build} succeeded.
The included compiler log records the kernel dependencies.
There are no \code{sorry} placeholders or user-declared
axioms in this checked target. Its reported dependencies
include \code{propext} and \code{Quot.sound}, and do not
include \code{Classical.choice}. This is a report about
the host proof terms; it is not a substitute for checking
the object theory in which a derivation is constructed.

The main checked declarations, in namespace
\code{ConstructiveIFT}, are:
\begin{center}
\small
\begin{tabular}{>{\raggedright\arraybackslash}p{0.42\textwidth}>{\raggedright\arraybackslash}p{0.50\textwidth}}
\toprule
Declaration & Content\\
\midrule
\code{errorBudget\_closed} &
Exact finite weighted geometric error sum.\\
\code{half\_contraction\_closed} &
The denominator-cleared inexact contraction inequality.\\
\code{half\_contraction\_budget\_finite} &
The bound using only the finite prefix of recurrence hypotheses.\\
\code{checked\_half\_contraction\_budget} &
A primitive-recursive-style Boolean trace checker implies
the finite bound.\\
\code{iterate\_add} &
Concatenation of two finite iteration segments.\\
\code{neumann\_residual} &
The finite vector identity \(z_N+E^Nb=Ez_N+b\), with
the required additive laws supplied explicitly.\\
\bottomrule
\end{tabular}
\end{center}

For example, the checked finite checker theorem has the form
\begin{lstlisting}
theorem checked_half_contraction_budget
    (d : Nat -> Nat) (e N : Nat)
    (checked : checkSteps d e N = true) :
    2 ^ N * d N <= errorBudget (d 0) e N
\end{lstlisting}
The host parameter \(d\) becomes a fixed primitive recursive
symbol in the object-language version. The checked theorem
does not by itself assert primitive recursiveness of an
arbitrary host function.

\subsection{The repository snapshot actually inspected}

The GitHub connector resolved ProveIt main to commit
\[
 \code{d36b35eed5bd01af4ade66bf539622495813ccc8}.
\]
The inspected sources and project status are recorded in
the archive's source manifest and audit. No repository
files were modified. The reference snapshot contains the
selected source modules, not their complete transitive
import closure, and those repository modules were not
rebuilt as part of this package.

The following paths are relative to
\path{Logic/PrimitiveRecursiveArithmetic/Derivations/Lean/IPRADerivations/}.
The source uses \(\Prf\,\PRAAx\), making it relevant to
object iPRA, rather than merely to analysis in host Lean.
\begin{center}
\small
\begin{tabular}{>{\raggedright\arraybackslash}p{0.35\textwidth}>{\raggedright\arraybackslash}p{0.57\textwidth}}
\toprule
Module & Inspected interface and reuse boundary\\
\midrule
\path{Foundations/ContractionTrace.lean} &
\code{contraction\_iter} uses an internal object index and
explicit \code{PRA.qfInd}. It is the best template for
the new additive-error trace invariant.\\[3pt]
\path{Analysis/HalfContractionTail.lean} &
Internal-index geometric and tail bounds, using a
nonincreasing finite potential.\\[3pt]
\path{Analysis/DyadicCauchy.lean} &
Closed object formulas for fixed primitive recursive
symbols with a particular adjacent dyadic-integer
motion hypothesis. An adapter to that format would
need proof; it is not automatic for our name.\\[3pt]
\path{Analysis/FiniteFTC.lean} &
Internal finite-mesh summation of signed local errors;
a suitable template for our polygon argument.\\[3pt]
\path{Foundations/ContractionIterates.lean} &
Internal exact-contraction iteration and Cauchy
templates. The rounded evaluator needs the inexact
adaptation.\\[3pt]
\path{Analysis/FeedbackInexact.lean} &
Genuine object proof terms, but the iteration count is
an external Lean index in the inspected theorem
family. It is not already the internally quantified
runtime loop theorem needed here.\\[3pt]
\path{Foundations/NeumannCorrection.lean} &
Genuine finite matrix identities in the object
theory, with an external truncation index.
An internal coded matrix loop is still required.\\
\bottomrule
\end{tabular}
\end{center}

These distinctions matter: an externally recursive family
of proof terms containing the numeral \(2^k\) is not the
same object theorem as a universally quantified bound
for variable \(k\). The existing modules supply valuable
proof patterns and elementary arithmetic, while the
new evaluator and the variable-length inexact argument
still require an integration layer.

\subsection{Required definitions in a full Lean development}

The detailed \code{lean/BLUEPRINT.md} accompanies this
section. A clean implementation has two connected levels.
Host Lean can use
\[
 \code{QVec n := Fin n -> Rat},\qquad
 \code{QMat m n := Fin m -> Fin n -> Rat},
\]
with explicit finite folds for norms. Object iPRA uses
their numerical list encodings. At either level the
necessary definitions are:
\begin{enumerate}
 \item Signed rational codes, their validity predicates,
 arithmetic, comparison, positive-denominator witnesses,
 and numerical definitions of finite vectors and matrices.
 \item Supremum and row-sum norms, product and triangle
 inequalities, boxes, dyadic floor, and clipping.
 \item Regular names, finite-observation equality and
 order, and the bounded-precision operations on names.
 \item The fixed evaluator symbols, raw certificates
 R0--R4, and the named-point extension programs with
 moduli \eqref{eq:rawenlarge}.
 \item The numerical trace function
 \eqref{eq:algorithm}, its clock \(N(k)\), and the finite
 prefix checker for local error inequalities.
 \item The primitive recursive graph evaluator,
 parameter-name composition, residual observation
 formulas, and quantitative uniqueness formulas.
 \item Coded matrix powers, finite Neumann sums, their
 rational perturbation bounds, and the derivative
 evaluator \eqref{eq:derivativeprogram}.
 \item The numerical Taylor and continuity output
 formulas with \(\mu_g,\nu_g\), and the rational
 preconditioner construction.
\end{enumerate}
All these definitions have been specified mathematically
above. None requires an unspecified function-space
quantifier in the arithmetic language.

\subsection{A concrete completion order}

\paragraph{1. Rational arithmetic and finite arrays.}
Implement the encodings and prove their primitive
recursive definitions satisfy the elementary field
identities under the denominator guards. Keep the
algorithmic rational operations shared between the
producer and the certificate checker. Prove clipping
nonexpansive by rational order cases.

\paragraph{2. The internally quantified inexact trace theorem.}
Port the integer invariant from the checked core to
\(\Prf\,\PRAAx\), using the internal-index method in
\code{ContractionTrace}. The induction formula is the
Boolean-prefix-check implication to the integer
invariant. This is the first main proof obligation
not already discharged by the supplied host-Lean file.

\paragraph{3. The raw Taylor-to-segment adapter.}
Implement \(L_p\), the finite rational polygon, and
the error budget at a fixed observation precision.
Use the finite signed-sum pattern in \code{FiniteFTC}.
Prove rational squeezing and instantiate it to obtain
the one-step contraction estimates for actual program
outputs. This step connects analytic input certificates
to the finite trace theorem.

\paragraph{4. The implicit-function name.}
Define the rounded trace by primitive recursion.
Use aligned final segments and the inexact trace
theorem to prove \eqref{eq:pairbound}; use adjacent
increments to prove \eqref{eq:residual}.
Build the regular output name directly. Then prove
the zero equation, uniqueness, residual bound,
Lipschitz estimate, and the parameter-name extension.
No separate theorem about an arbitrary complete
metric space is needed.

\paragraph{5. Internally variable finite matrix loops.}
Encode a matrix as a fixed-length numerical array.
Use componentwise rational loop invariants to prove
the two geometric identities and the power-difference
identity. The source in \code{NeumannCorrection}
provides exact finite algebra, but its external index
must be replaced by the coded internal loop. Prove
the finite norm sums \(2,4,8\) needed in
Lemma~\ref{lem:perturb}.

\paragraph{6. The \(C^1\) conclusion.}
Compose the Jacobian evaluator with the graph name at
the explicit index \(\nu(q+1)\), construct the
derivative name from a Neumann sum, and translate the
Taylor calculation and inverse-difference calculation
at arbitrary fixed output precision. This supplies
the exact formulas \(\mu_g,\nu_g\). Finally formalize
the determinant-apartness and translated-center
adapters.

\paragraph{7. Object-proof audit and an optional classical bridge.}
The final target should be an actual
\(\Prf\,\PRAAx\) object derivation of the finite
observation schema, with each induction constructor
carrying a verified quantifier-free formula.
Audit its dependency closure for placeholders and
unintended assumptions. A subsequent interpretation
in ordinary real analysis may use Mathlib's real
vectors and calculus to compare with its existing
implicit function theorems \cite{mathlib}.
Such a comparison is useful, but invoking Mathlib's
noncomputable implicit-function construction would
not establish the primitive recursive output or
the restricted-induction claim.

\subsection{Reproduction}

From the package root, run
\begin{lstlisting}
python3 code/validate.py
cd lean
lake build
\end{lstlisting}
The project pins Lean 4.32.0. To rebuild the report,
from the \code{paper} directory run
\begin{lstlisting}
pdflatex -interaction=nonstopmode -halt-on-error implicit_function_ipra.tex
pdflatex -interaction=nonstopmode -halt-on-error implicit_function_ipra.tex
\end{lstlisting}
The figure is already included. Its optional regeneration
script uses Matplotlib; Matplotlib is not required by
the rational solver or its validation suite.
